\paragraph{Proposition E.2.5 (Proved: the first integral torsion of the shifted product).}
Let \(Z_\varphi=X\times_B\varphi^*X\) be the transverse
fivefold in E.2.3, with the two sets of three critical values
disjoint. Its second integral homology is
\[
 H_2(Z_\varphi;\mathbb Z)
   =\mathbb Z\oplus(\mathbb Z/12)^7.
\]
In particular the rational Betti vector in E.2.3 cannot be
interpreted as integral torsion-free homology.
\begin{proof}
Use subscripts one and two for the two rank-four cohomology
local systems. At a critical value of one factor the other
factor is a smooth torus bundle over a disc, hence a product
there. Each local specialization is therefore the tensor
product of the original specialization with the identity on
the exterior algebra of \(\mathbb Z^4\). The original finite
reduced surfaces and the original cusp have free cohomology.
Their specialization maps are injective; at the cusp they
are integral isomorphisms onto invariants. For the latter,
apply the support-function construction of E.1.3 in two
coordinates. It supplies the invariant deck classes
\(b_1,b_2\), the three classes
\(a_1b_1,a_2b_2,a_1b_2+a_2b_1\), and their products.
Together with \(b_1b_2\) these give an integral kernel
basis in degree two; their products give integral bases
in degrees three and four. Both the central-fibre and
invariant lattices have ranks \((1,2,4,2,1)\), as follows
also from the normalization matrix in D.1.2 at \(d=1\).
Thus the specialization is onto between free groups of
equal ranks in every degree, proving the claimed isomorphism.
Thus the only difference between \(R^qg_*\mathbb Z\) and
\(\mathcal J_q=j_*\bigwedge^q(\mathcal V_1\oplus\mathcal V_2)\)
is their finite covering cokernel, with no point-supported kernel.

The finite Wang calculation (E.I1), now with the original
three-dimensional vertical lattice, gives
\[
 C_3^1=\mathbb Z/3,\quad C_4^1=\mathbb Z/4,
 \qquad C_3^2=0,\quad C_4^2=\mathbb Z/2.
\]
The invariant covector \(\gamma\) generates the first two
cokernels, and the invariant two-form \(q_X\) generates the
last, by the same even-intersection argument as in E.1.5.
Consequently the degree-two covering cokernel at all four
finite values, combined, is
\[
 C^2=\left((\mathbb Z/3)^4\oplus
                 (\mathbb Z/4)^4\oplus\mathbb Z/2\right)^2.
\]
The two monodromy groups act independently on the two factors.
The three global invariant two-forms have integral basis
\(q_1,\gamma_1\gamma_2,q_2\). The first and third map
onto the respective \(\mathbb Z/2\) summands and vanish
in the other summands. The middle form maps to the primitive
\(\gamma\) coordinate of each of the four order-three or
order-four blocks. After the two order-two summands have
been removed, it therefore generates a direct cyclic summand
of order twelve in \((\mathbb Z/12)^8\). Hence
\[
 \operatorname{coker}\{H^0(B,\mathcal J_2)\to C^2\}
      =(\mathbb Z/12)^7.
 \tag{E.20}
\]
The program \texttt{fourfolds\_shifted\_integral.py} realizes
this as an eighteen-row integer presentation; its Smith form
has eleven unit factors and seven factors twelve.

In degree two the global coboundary image is primitive.
Thus the parabolic group, a subgroup of the free-group
cohomology on the punctured base, is torsion free. Its rank
is eighteen: the six local invariant ranks are
\((16,16,18,16,16,18)\), and its global invariant and
coinvariant ranks are both three. The Euler formula on the
six-punctured sphere gives \(3+3+4\cdot28-100=18\).
The primitivity and all these ranks are checked by the exact
exterior matrices in the same program. The sheaf sequence
\(0\to R^2g_*\mathbb Z\to\mathcal J_2\to C^2\to0\)
and (E.20) now give
\[
 H^1(B,R^2g_*\mathbb Z)
     =\mathbb Z^{18}\oplus(\mathbb Z/12)^7.
\]
Here the quotient by its finite subgroup is free, so this
abelian-group extension splits.

The group \(H^2(B,R^1g_*\mathbb Z)\) is the free
rank-two group of coinvariants, one primitive summand from
each factor. Pulling back the integral degree-two
transgression of the original sphere fibration through each
projection kills these two summands separately and
surjectively. Thus \(E_\infty^{2,1}=0\) in the Leray
sequence of \(g\). Its other middle group
\(E_\infty^{0,3}\) is free, since \(R^3g_*\mathbb Z\)
injects into \(\mathcal J_3\). There are no higher
differentials. Consequently
\(\operatorname{Tor}H^3(Z_\varphi;\mathbb Z)
=(\mathbb Z/12)^7\). The universal coefficient theorem
identifies this with \(\operatorname{Tor}H_2(Z_\varphi;\mathbb Z)\).
Finally E.2.3 gives \(b_2=1\), proving the assertion.
This argument does not determine the extensions in the
higher integral Leray filtrations.
\end{proof}

